\subsection{Sensitivity Analysis}
\label{subsec:sensitivity}

To assess the robustness of our FDR framework to its two key hyperparameters---the number of negative samples $K$ and Storey's tuning parameter $\lambda$---and to the choice of FDR control procedure under dependence, we conducted a sensitivity analysis on the WN18RR benchmark.

\textbf{Sensitivity to $K$.}
Table~\ref{tab:sensitivity-K} reports $\hat{\pi}_0$ and BH rejection counts $R(0.05)$ for $K \in \{25, 50, 100\}$.
For TransE, $\hat{\pi}_0$ is identical at $K=50$ and $K=100$ (0.295), while RotatE shows modest variation (0.265 at $K=50$ vs.\ 0.259 at $K=100$, a difference of 0.006). BH rejection counts are sensitive at very small $K$ ($K=25$ produces $R=0$ for both models due to insufficient granularity of $p$-values at $1/26 \approx 0.038$ for $\alpha = 0.05$) but stabilize by $K=50$. RotatE's $R(0.05)$ increases from 505 ($K=50$) to 1,100 ($K=100$), reflecting the finer $p$-value resolution at larger $K$, while TransE plateaus at $K=50$. The recommended choice $K=100$ provides a conservative margin above the convergence threshold while remaining computationally accessible.

\begin{table}[h]
\centering
\caption{Sensitivity of FDR metrics to the number of negative samples $K$. $\lambda = 0.5$ fixed.}
\label{tab:sensitivity-K}
\begin{tabular}{lcccc}
\toprule
\textbf{Model} & \textbf{K=25} & \textbf{K=50} & \textbf{K=100} \\
\midrule
\multicolumn{4}{l}{\textit{$\hat{\pi}_0$ (Storey null proportion)}} \\
\quad TransE & 0.317 & 0.295 & 0.295 \\
\quad RotatE & 0.285 & 0.265 & 0.259 \\
\midrule
\multicolumn{4}{l}{\textit{$R(0.05)$ (BH rejections at $\alpha=0.05$)}} \\
\quad TransE & 0 (0.0\%) & 262 (9.0\%) & 262 (9.0\%) \\
\quad RotatE & 0 (0.0\%) & 505 (17.3\%) & 1,100 (37.6\%) \\
\bottomrule
\end{tabular}
\end{table}

\textbf{Sensitivity to $\lambda$.}
Table~\ref{tab:sensitivity-lambda} reports $\hat{\pi}_0$ for Storey's tuning parameter $\lambda \in \{0.3, 0.5, 0.7\}$.
The estimated $\hat{\pi}_0$ varies by at most 0.076 across this range, and the BH rejection count $R(0.05)$ is invariant to $\lambda$ (since BH does not depend on $\hat{\pi}_0$).
The $q$-value rejection counts vary from 1,843 to 2,050 (TransE) and 2,050 to 2,050 (RotatE) across $\lambda$ values---a range of approximately 10\%---which does not affect any of the paper's qualitative conclusions.

\begin{table}[h]
\centering
\caption{Sensitivity of $\hat{\pi}_0$ to Storey's tuning parameter $\lambda$. $K=100$ fixed.}
\label{tab:sensitivity-lambda}
\begin{tabular}{lcccc}
\toprule
\textbf{Model} & \textbf{$\lambda=0.3$} & \textbf{$\lambda=0.5$} & \textbf{$\lambda=0.7$} \\
\midrule
TransE & 0.330 & 0.295 & 0.270 \\
RotatE & 0.296 & 0.259 & 0.220 \\
\bottomrule
\end{tabular}
\end{table}

\textbf{Robustness to dependence structure: BY and adaptive BH procedures.}
The BH procedure guarantees FDR control under the PRDS condition (Section~\ref{subsec:bh-procedure}). To assess whether the core findings are robust to violations of this condition, we applied two alternative procedures that relax or eliminate the PRDS requirement.

\textbf{BY procedure.}
The Benjamini-Yekutieli (BY) procedure~\cite{benjamini2001control} provides FDR control under arbitrary dependence at the cost of reduced power.
The BY threshold is $\alpha / (m \cdot H_m)$ where $H_m = \sum_{i=1}^m 1/i$ is the $m$-th harmonic number; with $m = 2{,}924$, $H_m \approx 8.29$, making the per-hypothesis threshold approximately 165$\times$ more stringent than BH.
Consequently, BY produces zero rejections for all four models on both WN18RR ($R_{\text{BY}}(0.05) = 0$) and FB15k-237 at $\alpha = 0.05$---consistent with the known conservatism of BY when $m$ is large~\cite{benjamini2001control}.

\textbf{Two-stage adaptive BH.}
When $\pi_0 < 1$, the two-stage adaptive BH procedure~\cite{benjamini2006adaptive} recovers power by first estimating $m_0$ and then applying BH at an adjusted threshold $\alpha' = \alpha \cdot m / \hat{m}_0$.
This procedure retains finite-sample FDR control at level $\alpha$ under independence, without requiring the PRDS condition to be verified.
Table~\ref{tab:sensitivity-adapt} reports adaptive BH rejection counts alongside standard BH for all four models on WN18RR.
Adaptive BH substantially increases power (e.g., TransE: 262 $\to$ 1{,}505 at $\alpha = 0.05$; RotatE: 1{,}100 $\to$ 1{,}901) while preserving the same cross-model ordering (RotatE $>$ ConvE $\approx$ TransE $>$ ComplEx), confirming that the paper's comparative conclusions are not artifacts of the PRDS assumption.
ComplEx remains at zero rejections under all procedures, consistent with its $\hat{\pi}_0 = 0.754$.

\begin{table}[h]
\centering
\caption{Comparison of BH, BY, and two-stage adaptive BH rejection counts on WN18RR ($K=100$, $\alpha = 0.05$). Adaptive BH applies the Benjamini-Krieger-Yekutieli (2006) procedure, which estimates $m_0 = m \cdot \hat{\pi}_0$ and adjusts the BH threshold accordingly.}
\label{tab:sensitivity-adapt}
\begin{tabular}{lcccc}
\toprule
\textbf{Model} & \textbf{BH} & \textbf{BY} & \textbf{Adaptive BH} & $\boldsymbol{\hat{\pi}_0}$ \\
\midrule
TransE   & 262 (9.0\%)  & 0 (0.0\%) & 1{,}505 (51.5\%) & 0.295 \\
RotatE   & 1{,}100 (37.6\%) & 0 (0.0\%) & 1{,}901 (65.0\%) & 0.259 \\
ComplEx  & 0 (0.0\%)   & 0 (0.0\%) & 0 (0.0\%)       & 0.754 \\
ConvE    & 742 (25.4\%) & 0 (0.0\%) & 1{,}518 (51.9\%) & 0.321 \\
\bottomrule
\end{tabular}
\end{table}

These results indicate that the paper's core findings are robust to the choice of FDR control procedure: the existence of a non-trivial null proportion ($\hat{\pi}_0 \gg 0$) and the per-relation heterogeneity are established independently of the specific method, and the cross-model BH ordering (RotatE $>$ ConvE $\approx$ TransE $>$ ComplEx) is preserved under adaptive BH.
